% 3 ---------------------------------------------------------------------------
\begin{frame}{A Tree Partitions the Feature Space}
	\begin{columns}[c,onlytextwidth]
		\column{.47\textwidth}
		\centering
		\begin{tikzpicture}
			\node[decision] (r) at (0,0) {$x_1\le 5$?};
			\node[leaf] (a) at (-1.5,-1.4) {$R_1$\\output $w_1$};
			\node[decision] (b) at (1.5,-1.4) {$x_2\le 3$?};
			\node[leaf] (c) at (.35,-3) {$R_2$\\output $w_2$};
			\node[leaf] (d) at (2.65,-3) {$R_3$\\output $w_3$};
			\draw[treeline] (r)--node[above left,font=\scriptsize]{yes}(a);
			\draw[treeline] (r)--node[above right,font=\scriptsize]{no}(b);
			\draw[treeline] (b)--node[pos=.7,above left,font=\scriptsize]{yes}(c);
			\draw[treeline] (b)--node[pos=.7,above right,font=\scriptsize]{no}(d);
		\end{tikzpicture}
		\column{.47\textwidth}
		\centering
		\begin{tikzpicture}[x=.49cm,y=.49cm]
			\fill[bulletgreen!28] (0,0) rectangle (5,7);
			\fill[bulletgray!50] (5,0) rectangle (10,3);
			\fill[slideorange!12] (5,3) rectangle (10,7);
			\draw[flow] (0,0)--(10.6,0) node[right]{$x_1$};
			\draw[flow] (0,0)--(0,7.6) node[above]{$x_2$};
			\draw[thick] (5,0)--(5,7); \draw[thick] (5,3)--(10,3);
			\node at (2.5,3.5) {$R_1$}; \node at (7.5,1.5) {$R_2$}; \node at (7.5,5) {$R_3$};
			\node[below] at (5,0) {5}; \node[left] at (0,3) {3};
		\end{tikzpicture}
	\end{columns}
	\pause
	\vspace{.3cm}
	\[f(x)=\sum_{v=1}^{J}w_v\,\mathbf 1\{x\in R_v\}.\]
	\centering Each leaf defines a region with a constant prediction.
\end{frame}

% 4 ---------------------------------------------------------------------------
\begin{frame}{What Does a Leaf Predict?}
	\begin{columns}[T,onlytextwidth]
		\column{.47\textwidth}
		{\large\color{slideorange}\textbf{Regression}}
		\vspace{.25cm}
		\par For squared error, use the mean target in the leaf:
		\[w_v=\frac{1}{|I_v|}\sum_{i\in I_v}y_i.\]
		\vspace{.15cm}
		\centering
		\begin{tikzpicture}
			\node[leaf,text width=4.2cm] {Targets: $2,4,6$\\[4pt]Prediction: $4$};
		\end{tikzpicture}
		\column{.47\textwidth}
		{\large\color{slideorange}\textbf{Classification}}
		\vspace{.25cm}
		\par Estimate class proportions in the leaf:
		\[p_v(c)=\frac{\sum_{i\in I_v}\mathbf 1\{y_i=c\}}{|I_v|}.\]
		\vspace{.15cm}
		\centering
		\begin{tikzpicture}
			\node[leaf,text width=4.2cm] {Labels: $0,1,1,1$\\[4pt]$p_v(1)=0.75$; predict class $1$};
		\end{tikzpicture}
	\end{columns}
	\vspace{.6cm}
	$I_v$: training samples reaching leaf $v$; $J$: number of leaves.
\end{frame}

% 5 ---------------------------------------------------------------------------
\begin{frame}{Choose a Criterion, Then Compare Splits}
	Split the samples $I$ into children $I_L$ and $I_R$.\par
	Let $A(I)$ be node impurity or average prediction loss.
	\[
	\Delta=A(I)-\frac{|I_L|}{|I|}A(I_L)-\frac{|I_R|}{|I|}A(I_R).
	\]
	\centering Select the allowed split with the largest decrease.
	\pause
	\vspace{.45cm}
	\begin{columns}[T,onlytextwidth]
		\column{.47\textwidth}
		{\large\color{slideorange}\textbf{Regression}}\par
		\vspace{.2cm}
		Squared error: use a leaf mean.\par
		\vspace{.2cm}
		Absolute error: use a leaf median.
		\column{.47\textwidth}
		{\large\color{slideorange}\textbf{Classification}}\par
		\vspace{.2cm}
		Gini impurity or entropy.\par
		\vspace{.2cm}
		Prefer purer child nodes.
	\end{columns}
	\src{Weight each child by its fraction of the samples. The criterion is a modeling choice.}
\end{frame}

% 6 ---------------------------------------------------------------------------
\begin{frame}{Regression: Minimize Squared Error}
	\setlength{\abovedisplayskip}{5pt}\setlength{\belowdisplayskip}{5pt}
	\setlength{\abovedisplayshortskip}{3pt}\setlength{\belowdisplayshortskip}{5pt}
	\textbf{SSE: sum of squared errors.} For a leaf containing samples $I$,
	\[\mathrm{SSE}(I)=\sum_{i\in I}(y_i-\bar y_I)^2,
	\qquad \mathrm{MSE}(I)=\frac{\mathrm{SSE}(I)}{|I|}.\]
	\centering
	\renewcommand{\arraystretch}{1.1}
	\begin{tabular}{c|rrrr}
		Feature $x$ & 1 & 2 & 3 & 4\\
		Target $y$ & 1 & 2 & 8 & 9
	\end{tabular}
	\vspace{.5cm}
	\begin{columns}[T,onlytextwidth]
		\column{.46\textwidth}
		\textbf{One leaf}
		\[\bar y=5,\qquad \mathrm{SSE}=50.\]
		\column{.49\textwidth}
		\pause
		\textbf{Split at $x\le2.5$}
		\[w_L=1.5,\qquad w_R=8.5,\]
		\[\mathrm{SSE}_L+\mathrm{SSE}_R=1.\]
	\end{columns}
	\pause
	\vspace{.05cm}
	\[\Delta_{\mathrm{MSE}}=
	\frac{\mathrm{SSE}(I)-\mathrm{SSE}(I_L)-\mathrm{SSE}(I_R)}{|I|}.\]
\end{frame}

% 7 ---------------------------------------------------------------------------
\begin{frame}{Classification: Gini or Entropy}
	\setlength{\abovedisplayskip}{5pt}\setlength{\belowdisplayskip}{5pt}
	\setlength{\abovedisplayshortskip}{3pt}\setlength{\belowdisplayshortskip}{5pt}
	Let $p_c$ be the fraction of samples in class $c$ at the node.
	\vspace{.3cm}
	\begin{columns}[T,onlytextwidth]
		\column{.47\textwidth}
		{\large\color{slideorange}\textbf{Gini impurity}}\par
		\[G(I)=1-\sum_{c=1}^{C}p_c^2.\]
		Measures class mixing.
		\column{.47\textwidth}
		\pause
		{\large\color{slideorange}\textbf{Entropy}}\par
		\[H(I)=-\sum_{c=1}^{C}p_c\log_2 p_c.\]
		Label uncertainty in bits; $0\log_2 0=0$.
	\end{columns}
	\pause
	\vspace{.2cm}
	\centering
	\renewcommand{\arraystretch}{1.15}
	\begin{tabular}{lcc}
		\toprule
		Binary class proportions & Gini & Entropy\\ \midrule
		Pure: $(1,0)$ & $0$ & $0$\\
		Balanced: $(0.5,0.5)$ & $0.5$ & $1$\\ \bottomrule
	\end{tabular}
	\vspace{.2cm}
	\par Both favor pure nodes, but may choose different splits.
\end{frame}

% 8 ---------------------------------------------------------------------------
\begin{frame}{One Split, Two Impurity Measures}
	\setlength{\abovedisplayskip}{5pt}\setlength{\belowdisplayskip}{5pt}
	\setlength{\abovedisplayshortskip}{3pt}\setlength{\belowdisplayshortskip}{5pt}
	\centering
	\renewcommand{\arraystretch}{1.1}
	\begin{tabular}{lcccc}
		\toprule
		Node & Class 0 & Class 1 & Gini $G$ & Entropy $H$\\ \midrule
		Parent & 4 & 4 & $0.500$ & $1.0000$\\
		Left child & 3 & 1 & $0.375$ & $0.8113$\\
		Right child & 1 & 3 & $0.375$ & $0.8113$\\ \bottomrule
	\end{tabular}
	\pause
	\vspace{.1cm}
	\par\textbf{Gini decrease}
	\[\Delta_G=0.5-\frac48(0.375)-\frac48(0.375)=0.125.\]
	\pause
	\par\textbf{Entropy decrease: information gain}
	\[\Delta_H=1-\frac48(0.8113)-\frac48(0.8113)\approx0.1887\text{ bits}.\]
	\centering Gain scales differ: compare candidates within one criterion.
\end{frame}

% 9 ---------------------------------------------------------------------------
\begin{frame}{Grow the Tree One Split at a Time}
	\begin{enumerate}\spaceditems
		\item Start with all training samples at the root.
		\item At a node, score candidate feature--threshold pairs using the chosen criterion.
		\pause
		\item Choose the best allowed split and send samples to its children.
		\item Repeat until a stopping condition is reached; assign leaf outputs.
	\end{enumerate}
	\pause
	\vspace{.35cm}
	\begin{columns}[T,onlytextwidth]
		\column{.47\textwidth}
		\textbf{Local search}\par
		An early split changes every later decision in its subtree.
		\column{.47\textwidth}
		\textbf{Computational shortcut}\par
		Sort a numerical feature and scan candidate cut points.
	\end{columns}
	\vspace{.35cm}
	\centering Greedy splitting does not guarantee the globally best tree.
\end{frame}

% 10 ---------------------------------------------------------------------------
\begin{frame}{Control How Much the Tree Can Fit}
	\begin{columns}[c,onlytextwidth]
		\column{.46\textwidth}
		\centering
		\begin{tikzpicture}[scale=.85]
			\node[decision] (r) at (0,0) {split};
			\node[leaf] (a) at (-1.2,-1.4) {mean};
			\node[leaf] (b) at (1.2,-1.4) {mean};
			\draw[treeline] (r)--(a); \draw[treeline] (r)--(b);
			\node[align=center] at (0,-2.4) {Few broad regions};
		\end{tikzpicture}
		\column{.5\textwidth}
		\begin{itemize}\spaceditems
			\item Limit maximum depth.
			\item Require enough samples in each leaf.
			\item Prune a larger tree using a complexity penalty.
		\end{itemize}
	\end{columns}
	\pause
	\vspace{.25cm}
	\[\text{Pruning objective:}\qquad R(T)+\alpha\,|\operatorname{leaves}(T)|.\]
	\centering $R(T)$: total training loss or impurity across leaves.\par
	\vspace{.2cm}
	Choose complexity on validation data; deeper trees can have higher variance.
\end{frame}

% 11 ---------------------------------------------------------------------------
\begin{frame}{Learn Where Missing Values Should Go}
	Using Gini impurity, compare both missing-value directions for $x\le5$.
	\vspace{.3cm}
	\centering
	\renewcommand{\arraystretch}{1.15}
	\begin{tabular}{c|rrrrrr}
		$x$ & 1 & 2 & 8 & 9 & missing & missing\\
		$y$ & 0 & 0 & 1 & 1 & 1 & 1
	\end{tabular}
	\pause
	\vspace{.35cm}
	\begin{tabular}{lccc}
		\toprule
		Missing direction & Left labels & Right labels & Weighted Gini\\ \midrule
		Left & $0,0,1,1$ & $1,1$ & $1/3$\\
		\textcolor{forestgreen}{Right} & $0,0$ & $1,1,1,1$ & \textcolor{forestgreen}{$0$}\\ \bottomrule
	\end{tabular}
	\pause
	\vspace{.35cm}
	\par Learn $(\text{feature},\text{threshold},\text{missing direction})$ at each node.\par
	\vspace{.2cm}
	This is a routing rule. Later splits can still separate missing samples.
\end{frame}

% 12 ---------------------------------------------------------------------------
\begin{frame}{A Surrogate Split Uses Another Feature}
	\begin{columns}[c,onlytextwidth]
		\column{.48\textwidth}
		\centering
		\begin{tikzpicture}[node distance=.65cm]
			\node[decision,text width=4.5cm] (a) {Primary rule\\income $\le a$?};
			\node[orangebox,text width=4.5cm,below=of a] (b) {If income is missing\\use assets $\le b$?};
			\node[leaf,below left=.65cm and -.4cm of b] (l) {left};
			\node[leaf,below right=.65cm and -.4cm of b] (r) {right};
			\draw[flow] (a)--(b); \draw[flow] (b)--(l); \draw[flow] (b)--(r);
		\end{tikzpicture}
		\column{.48\textwidth}
		\pause
		\begin{itemize}\spaceditems
			\item Find a rule that reproduces the primary split's left/right assignments.
			\item Samples missing the primary feature may take different branches.
			\item A learned default direction can also separate samples later, using other features.
		\end{itemize}
	\end{columns}
	\src{CART-style surrogate splits; illustrative rules.}
\end{frame}

% 13 ---------------------------------------------------------------------------
\begin{frame}{A Small Data Change Can Change the Tree}
	\centering
	\begin{tikzpicture}
		\node[decision] (a) at (-3.3,0) {$x_1\le a$?};
		\node[leaf] (al) at (-4.5,-1.4) {subtree 1};
		\node[leaf] (ar) at (-2.1,-1.4) {subtree 2};
		\node[decision] (b) at (3.3,0) {$x_2\le b$?};
		\node[leaf] (bl) at (2.1,-1.4) {subtree 3};
		\node[leaf] (br) at (4.5,-1.4) {subtree 4};
		\draw[treeline] (a)--(al); \draw[treeline] (a)--(ar);
		\draw[treeline] (b)--(bl); \draw[treeline] (b)--(br);
		\draw[flow] (-1.9,0)--node[above,font=\small,align=center]{perturb the\\training sample}(1.9,0);
	\end{tikzpicture}
	\pause
	\vspace{.55cm}
	\begin{itemize}\spaceditems
		\item Competing splits may have very similar training scores.
		\item A different early split changes the downstream partition.
		\item Averaging many diverse trees can stabilize predictions.
	\end{itemize}
	\src{Schematic illustration of tree instability. Next: random forests.}
\end{frame}
